\documentclass[12pt]{article}
\usepackage[utf8]{inputenc}
\usepackage[english]{babel}
\usepackage{geometry}
\geometry{a4paper, margin=2.5cm}
\usepackage{graphicx}
\usepackage{hyperref}

\title{Demonstration, Validation and Tests}
\author{NeoLAAC Development Team}
\date{\today}

\begin{document}

\maketitle

\section{General Introduction and Demonstration Methodology}
The validation and demonstration phase of the NeoLAAC (Freshman Support Line) system focuses on verifying the functional correctness, responsiveness, and visual integrity of the decentralized microservices ecosystem. The primary objective is to demonstrate that the user interface, routing mechanisms, database transactions, and emergency response workflows operate seamlessly under simulated real-world conditions. This document maps each of the implementation modules to specific visual validation checkpoints and user interface screenshots. Rather than exploring the source code repository directly, this report validates system performance through graphical outputs, end-to-end interface execution, and automated test reports.

\section{Central Orchestrator and API Gateway}
This section presents the validation of the system's ingress data layer and routing. The central orchestrator manages external traffic and acts as the gatekeeper for all internal communication, translating client actions into backend calls, as demonstrated in the following visual validations.

\subsection{orq Service}
The validation of the orq service focuses on verifying that the API gateway correctly routes incoming HTTP queries to the respective services and acts on critical messages. The graphical interface validation shows the chat window intercepting messages containing emergency keywords, triggering the responder warning system.

\begin{figure}[htbp]
    \centering
    % 
\includegraphics[width=0.8\textwidth]{images/orq_chat_triage.png}
    \caption{User interface demonstration of the orq service routing engine intercepting a simulated student chat message containing the critical keyword emergency, triggering a fallback response and dispatching an alert to the staff dashboard.}
    \label{fig:orq_chat_triage}
\end{figure}

\section{Core User Services and Access Control}
This section covers the visual validation of credentials, profiles, audit trails, and student session tracking. These interfaces ensure that access to student portals is authenticated and that logs are generated correctly for system audits.

\subsection{auth-service Service}
The auth-service is validated via the user login panel. The interface demonstrates successful credentials validation and JWT token issuance. The login interface is verified to exclude the previously implemented captcha challenge, since local dynamic IP allocations caused external captcha validation queries to fail.

\begin{figure}[htbp]
    \centering
    % \includegraphics[width=0.7\textwidth]{images/login_panel.png}
    \caption{Visual demonstration of the login interface showing a successful login response, highlighting the simplified credentials panel where the captcha box was removed to ensure continuous access over dynamic IP networks.}
    \label{fig:login_panel}
\end{figure}

\subsection{profile-service Service}
The profile-service is demonstrated through the user profile page. The UI displays editable display names, contact phone numbers, and profile avatars, validating that students can manage their identifiers independently.

\begin{figure}[htbp]
    \centering
    % \includegraphics[width=0.8\textwidth]{images/profile_page.png}
    \caption{Demonstration of the student profile management screen showing display name customization and avatar uploads successfully rendered in the browser.}
    \label{fig:profile_page}
\end{figure}

\subsection{logU-service Service}
The logU-service is verified by auditing user-related events in the administrative console. The validation displays specific audit messages, such as successful authentication events and changes to privilege roles, confirming correct logging.

\begin{figure}[htbp]
    \centering
    % \includegraphics[width=0.9\textwidth]{images/user_audit_logs.png}
    \caption{Administrative console interface showing user access history logs generated during student register and login operations.}
    \label{fig:user_audit_logs}
\end{figure}

\subsection{Usertrack-service Service}
The Usertrack-service is demonstrated in the responder dashboard through the volunteer status tracker. The interface displays active support agents with green indicators for online presence, confirming the routing algorithm can dispatch chats correctly.

\begin{figure}[htbp]
    \centering
    % \includegraphics[width=0.8\textwidth]{images/volunteer_tracking.png}
    \caption{Staff dashboard view illustrating the real-time presence indicators of volunteers and support personnel logged into the system.}
    \label{fig:volunteer_tracking}
\end{figure}

\section{Academic and Information Services}
This section validates the curriculum catalog, calendar schedules, automatic date generation, and informational pages. These screens ensure that students receive accurate curriculum timelines.

\subsection{academic-service Service}
The academic-service is demonstrated through the course directory page. The UI lists official courses, curricular years, and subjects query responses, verifying that the central curriculum data is accessible to students.

\begin{figure}[htbp]
    \centering
    % \includegraphics[width=0.8\textwidth]{images/courses_catalog.png}
    \caption{Freshman portal view displaying the official academic course catalog and subject details retrieved from the academic database.}
    \label{fig:courses_catalog}
\end{figure}

\subsection{calendar-service Service}
The calendar-service is validated by rendering the student academic calendar. The UI illustrates a month-by-month grid mapping class slots and scheduled evaluation dates.

\begin{figure}[htbp]
    \centering
    % \includegraphics[width=0.8\textwidth]{images/calendar_view.png}
    \caption{Calendar page interface illustrating the semester schedule grid displaying classes and academic activities for the selected course.}
    \label{fig:calendar_view}
\end{figure}

\subsection{ClassDates Integration}
The ClassDates library integration is validated by generating class occurrences. The interface displays a calculated list of exact class dates for a semester, confirming that pauses like Christmas break and national holidays were excluded from the generation.

\begin{figure}[htbp]
    \centering
    % \includegraphics[width=0.8\textwidth]{images/classdates_generation.png}
    \caption{Administrative interface showing the automatic generation of class schedules using the ClassDates library developed by Professor Dr. Tiago M. C. Simões.}
    \label{fig:classdates_generation}
\end{figure}

\subsection{FAQ-service Service}
The FAQ-service is validated via the help center interface. The screen displays a categorization of questions and answers regarding dining, Wi-Fi networks, and library schedules.

\begin{figure}[htbp]
    \centering
    % \includegraphics[width=0.8\textwidth]{images/faq_list.png}
    \caption{Help center interface demonstrating structured FAQ items grouped by student categories such as General, Technology, and Dining.}
    \label{fig:faq_list}
\end{figure}

\subsection{chatbot-service Service}
The chatbot-service is demonstrated via the virtual mentor chat window. The UI displays the chatbot matching a normalized student query against the FAQ data and returning an answer.

\begin{figure}[htbp]
    \centering
    % 
\includegraphics[width=0.8\textwidth]{images/chatbot_chat.png}
    \caption{Interactive chat interface showing the digital mentor matching normalized queries to provide answers regarding campus services.}
    \label{fig:chatbot_chat}
\end{figure}

\subsection{page-service Service}
The page-service is validated through static guides pages. The interface displays UBI departments, emergency contacts, local bus schedules, and partner discounts.

\begin{figure}[htbp]
    \centering
    % \includegraphics[width=0.8\textwidth]{images/info_guides.png}
    \caption{Survival guides portal displaying department contact directories and general integration tips for incoming students.}
    \label{fig:info_guides}
\end{figure}

\section{Social Interaction and Gamification Services}
This section describes the validation of forum posts, community interactions, student organizations, and gamified challenges. These interfaces confirm that freshmen can build communities and explore the campus.

\subsection{challenge-service Service}
The challenge-service is demonstrated via the active competitions panel. The UI displays active challenges, the photo upload field, captions inputs, and the gallery of approved student submissions.

\begin{figure}[htbp]
    \centering
    % \includegraphics[width=0.8\textwidth]{images/challenges_ui.png}
    \caption{Photo challenge screen showing an active competition detail page, the photo upload zone, and the gallery of student submissions.}
    \label{fig:challenges_ui}
\end{figure}

\subsection{post-service Service}
The post-service is validated by publishing student posts. The UI shows a community board where students create posts with media attachments, click the like icon, and submit comments.

\begin{figure}[htbp]
    \centering
    % \includegraphics[width=0.8\textwidth]{images/social_feed.png}
    \caption{Community forum board displaying user posts containing images, calculated likes count, and active comment sections.}
    \label{fig:social_feed}
\end{figure}

\subsection{news-service Service}
The news-service is demonstrated in the home dashboard news feed. The interface displays formal institution announcements, confirming a separate channel for verified academic news.

\begin{figure}[htbp]
    \centering
    % 
\includegraphics[width=0.8\textwidth]{images/news_feed.png}
    \caption{Freshman portal home dashboard displaying official academic news and announcements separated from user posts.}
    \label{fig:news_feed}
\end{figure}

\subsection{events-service Service}
The events-service is validated by displaying student organization event schedules. The interface renders cards for upcoming workshops, sports events, and academic integration meetings.

\begin{figure}[htbp]
    \centering
    % \includegraphics[width=0.8\textwidth]{images/events_cards.png}
    \caption{Student association events feed showing detailed cards for upcoming recreational and integration activities.}
    \label{fig:events_cards}
\end{figure}

\subsection{feed-service Service}
The feed-service is demonstrated on the social home feed, verifying that public posts, announcements, and events are aggregated into a single chronological stream.

\begin{figure}[htbp]
    \centering
    % \includegraphics[width=0.8\textwidth]{images/unified_feed.png}
    \caption{Unified portal timeline showcasing the aggregation of posts, events, and news in chronological order.}
    \label{fig:unified_feed}
\end{figure}

\subsection{comunidade-service Service}
The comunidade-service is validated through the student directory. The interface shows followers lists, student profiles, and options to follow other peers.

\begin{figure}[htbp]
    \centering
    % \includegraphics[width=0.8\textwidth]{images/student_networking.png}
    \caption{Student directory screen displaying follower list counters and networking interactions on student profiles.}
    \label{fig:student_networking}
\end{figure}

\section{Service of Communication, Support and Emergency Response}
This section details the validation of the communication tools, silent distress buttons, urgent phone calls, ticket trackers, and notification alerts. These interfaces confirm that freshmen can reach support staff.

\subsection{chat-service Service}
The chat-service is validated by initiating a chat support conversation. The interface shows the chat thread between a student and a support agent, demonstrating message delivery and historical chat logging.

\begin{figure}[htbp]
    \centering
    % 
\includegraphics[width=0.8\textwidth]{images/chat_thread.png}
    \caption{Active chat thread window displaying messages exchanged between a student in need and a support responder.}
    \label{fig:chat_thread}
\end{figure}

\subsection{emergency-service Service}
The emergency-service is validated via the distress button panel. The screen displays the panic button interface that, when triggered, sends latitude and longitude geolocalizations to the dispatch console.

\begin{figure}[htbp]
    \centering
    % 
\includegraphics[width=0.7\textwidth]{images/emergency_trigger.png}
    \caption{Mobile interface of the student app illustrating the panic button used to trigger emergency alerts and send GPS coordinates.}
    \label{fig:emergency_trigger}
\end{figure}

\subsection{emergencyCall-service Service}
The emergencyCall-service is validated via the callback request panel. The interface displays the request confirmation showing the student's phone number and the dispatcher phone status.

\begin{figure}[htbp]
    \centering
    % 
\includegraphics[width=0.8\textwidth]{images/call_request.png}
    \caption{Callback request interface displaying the phone connection status between the student and the responder dispatch.}
    \label{fig:call_request}
\end{figure}

\subsection{Emergency Services Separation Rationale}
The architectural separation of emergency, chat, and call services ensures that high-priority distress alerts remain unaffected by other application traffic. A visualization of this architecture is displayed below, illustrating the routing channels.

\begin{figure}[htbp]
    \centering
    % \includegraphics[width=0.8\textwidth]{images/emergency_architecture.png}
    \caption{Architectural flow map demonstrating the logical separation of low-latency emergency alerts, real-time chat logging, and telephony integrations.}
    \label{fig:emergency_architecture}
\end{figure}

\subsection{ticket-service Service}
The ticket-service is demonstrated via the support dashboard. The interface shows academic tickets queue status, team assignments, and resolution notes inputs.

\begin{figure}[htbp]
    \centering
    % \includegraphics[width=0.8\textwidth]{images/tickets_dashboard.png}
    \caption{Support team dashboard displaying the list of active tickets, priority levels, and assigned personnel.}
    \label{fig:tickets_dashboard}
\end{figure}

\subsection{notification-service Service}
The notification-service is validated by triggering a push notification. The frontend interface displays a real-time notification banner confirming a new message or alert has arrived.

\begin{figure}[htbp]
    \centering
    % \includegraphics[width=0.8\textwidth]{images/push_notification.png}
    \caption{User interface popup demonstrating a real-time push notification banner alerting staff of a pending emergency.}
    \label{fig:push_notification}
\end{figure}

\section{Geolocation, Maps and Routing}
This section covers the visual validation of geographic maps, campus buildings directories, and route computation. These maps verify that freshmen can find campuses, block rooms, and bus stations.

\subsection{map-service Service}
The map-service is validated through the campus map panel. The Leaflet.js interface shows UBI campuses markers and computes routes from student coordinates, displaying walking directions or bus stop routes.

\begin{figure}[htbp]
    \centering
    % \includegraphics[width=0.8\textwidth]{images/map_routing.png}
    \caption{Map navigation panel displaying the calculated route path on Leaflet.js and walking instructions toward the campus.}
    \label{fig:map_routing}
\end{figure}

\subsection{map-emergency-service Service}
The map-emergency-service is demonstrated on the responder dashboard map. The Leaflet.js interface displays active emergency pins geolocalized in real-time, mapping the responder's path toward the student.

\begin{figure}[htbp]
    \centering
    % \includegraphics[width=0.8\textwidth]{images/responder_map.png}
    \caption{Responder console map plotting active emergency location markers alongside computed rescue routes.}
    \label{fig:responder_map}
\end{figure}

\section{Database and Network Infrastructure}
This section validates relational schemas, database dumps, Cloudflare Zero Trust gateways, SSL certificates, Nginx static servers, and compose parameters. These confirm infrastructure stability.

\subsection{DB-central Service}
The DB-central service database is validated via the MariaDB console, showing correct execution of schema tables including users, logs, and tickets.

\begin{figure}[htbp]
    \centering
    % \includegraphics[width=0.9\textwidth]{images/database_console.png}
    \caption{Terminal interface showing database console displaying the active relational tables in the central MariaDB container.}
    \label{fig:database_console}
\end{figure}

\subsection{backup-service Service}
The backup-service is validated by inspecting the persistent backup directory. The interface lists timestamped SQL files containing dumps of the relational databases.

\begin{figure}[htbp]
    \centering
    % \includegraphics[width=0.8\textwidth]{images/backups_list.png}
    \caption{File browser view showing the timestamped backup SQL files generated automatically by the backup service.}
    \label{fig:backups_list}
\end{figure}

\subsection{tunnel-service Service}
The tunnel-service is validated in the Cloudflare Zero Trust dashboard. The interface shows the tunnel status as active, confirming outbound traffic routing without open port dependencies.

\begin{figure}[htbp]
    \centering
    % \includegraphics[width=0.9\textwidth]{images/cloudflare_tunnel.png}
    \caption{Cloudflare dashboard interface illustrating the active connection status and route configurations of the Zero Trust tunnel.}
    \label{fig:cloudflare_tunnel}
\end{figure}

\subsection{certbot Service}
The certbot service is validated by checking the Nginx security configuration. The output shows Let's Encrypt keys mapped, confirming automated ACME HTTP-01 challenge validations.

\begin{figure}[htbp]
    \centering
    % \includegraphics[width=0.8\textwidth]{images/ssl_status.png}
    \caption{Terminal log output displaying successful ACME challenges and Let's Encrypt SSL certificates renewal records.}
    \label{fig:ssl_status}
\end{figure}

\subsection{frontend-service Service}
The frontend-service is validated by loading the portal homepage. The UI shows Nginx successfully serving static HTML pages on port 8080.

\begin{figure}[htbp]
    \centering
    % 
\includegraphics[width=0.8\textwidth]{images/frontend_home.png}
    \caption{Web browser interface displaying the NeoLAAC portal homepage served by Nginx.}
    \label{fig:frontend_home}
\end{figure}

\subsection{dev-service Service}
The dev-service setup is validated via the Docker Compose terminal output. The log shows container initializations, network setups, and active environment flags.

\begin{figure}[htbp]
    \centering
    % \includegraphics[width=0.8\textwidth]{images/docker_compose.png}
    \caption{Terminal log output showing successful execution of docker-compose up starting all services.}
    \label{fig:docker_compose}
\end{figure}

\subsection{logging-service Service}
The logging-service is validated via the admin logging dashboard, illustrating centralized operational logs from authentication and emergency services.

\begin{figure}[htbp]
    \centering
    % \includegraphics[width=0.8\textwidth]{images/logging_dashboard.png}
    \caption{Administrative log viewer illustrating aggregated debug and warning events from all services.}
    \label{fig:logging_dashboard}
\end{figure}

\section{Quality Control and Testing Pipeline}
This section validates the automatic testing pipeline designed to prevent code regression. The pipeline executes test blocks in an isolated environment, ensuring continuous integration, as detailed below.

\subsection{unit-tests Service}
The unit-tests execution is validated via the test logs. The output displays passing assertions for api routing and token validation, proving code stability.

\begin{figure}[htbp]
    \centering
    % \includegraphics[width=0.8\textwidth]{images/tests_execution.png}
    \caption{Terminal logs showing passing pytest assertions for authentication and gateway routes.}
    \label{fig:tests_execution}
\end{figure}

\subsection{pipeline.py Orchestration}
The pipeline.py execution is validated via the generated report. The interface shows test_report.md matching all green status, verifying that containers were torn down successfully after execution.

\begin{figure}[htbp]
    \centering
    % \includegraphics[width=0.8\textwidth]{images/pipeline_report.png}
    \caption{File browser view showing the generated test report markdown file indicating all checks passed.}
    \label{fig:pipeline_report}
\end{figure}

\section{System Verification and Test Execution Results}
The NeoLAAC quality control strategy is validated through comprehensive unit and integration test suites. The execution results verify that the individual microservices communicate reliably and that critical APIs enforce authorization rules.

The test suite in the authentication service validates ten separate test cases. These tests assert user registration behaviors, validating that attempts to register duplicate email addresses return a bad request status. Authentication routes verify admin credential access and normal student login, asserting that correct passwords return JSON Web Tokens and incorrect passwords result in unauthorized error codes. Role management and token validation tests assert that decodable JWT signatures return true sub fields, and invalid tokens are rejected. Database commits are mocked to ensure tests run in isolation.

The orchestrator integration test suite verifies over twenty API routes. The gateway tests check database health indicators and user listing routes, asserting correct data proxying. Chat messaging tests simulate multiple flows, asserting that messages containing critical keywords (such as emergency, panic, or accident) are marked as escalated and dispatch real-time notifications to the response team role. Chatbot fallback scenarios are tested by sending unresolvable queries, asserting that the orchestrator escalates the session to human staff when the chatbot returns a comprehension warning. Geolocation navigation tests simulate route queries, asserting that the map service returns OSRM pedestrian routing paths.

The execution report generated by the smart testing pipeline confirms that all build and test steps passed successfully. Contentor images for auth-service, ticket-service, orchestrator, emergency-service, and profile-service compiled without errors. The test execution logs show all assertions resolved successfully, with a final status of all green, validating the robustness of the implementation.

\end{document}
